\section{The positive small-field measure and its complement}
\label{subsec:interacting-complement}

\cite[Proposition~3.1]{TerminalPressure} identifies $F_N$
with $(A_*;\zeta_j)_{j<k_*}$. Their scalar
coordinates have the form
$\zeta_j(\Upsilon)=\langle\phi_{j,\Upsilon},F^{\rm u}\rangle$
with orthonormal $\phi_{j,\Upsilon}$; each
coarse plaquette coordinate $(dA_*)(p)$ has
the form $\langle\phi_p,F^{\rm u}\rangle$ with
$\|\phi_p\|\le1$.
There are at most $2\ell_*^3L^{3(k_*-j)}$
coordinates at level $j$ and $3\ell_*^3$
coarse plaquettes. The stop Gaussian
$\dd\mathsf m_*$ is nondegenerate on the
real mean-zero coclosed coarse fields.
The row and count statements are
\cite[(13)]{TerminalPressure} and
\cite[(14)]{TerminalPressure}; the
Gaussian assertion is \cite[(12)]{TerminalPressure}.

Let $\mathfrak G$ require
$|\zeta_j(\Upsilon)|\le p_{0,j}$
and $|dA_*|_\infty\le p_*$, and let
$\mathfrak G^-$ decrease each threshold by one.
Put
\begin{align}
 \varpi_0&=18\ell_*^3\mathrm e^{-p_{0*}/2}
               +6\ell_*^3\mathrm e^{-(p_*-1)^2/(2C_0^2)},\notag\\
 t_2&=e\,\min\left\{N_*^{-1/2},\,C_\beta^{-1},
                                     (2\sqrt{N_*\mathbf D})^{-1}\right\},\notag\\
 \varpi_L&=\sqrt{\varpi_0}\,\mathrm e^{\mathbf D/2},
 \qquad M_L=2+\mathrm e^{\mathbf D/2},\notag\\
 \lambda_L&=(\varpi_L^{1/3}M_L^{2/3})^{1/N_*}
                              5(2N_*r/t_2)^4 .
                                     \label{eq:interacting-complement-constants}
\end{align}
A reciprocal with zero denominator is interpreted
as $+\infty$. The number $\mathbf D$ is the
uniform probe energy bound, so $v_\eta\le\mathbf D$.

\begin{proposition}\label{prop:interacting-complement}
Assume the hypotheses of \cref{prop:interacting-one-cell,prop:interacting-mixed-logarithm}
and the tail conditions $v_*\ge8$ and $p_0v_*^{p_0-1}\ge28$.
For a fixed $0<s<1$, set
$\bar\kappa_1=6e_*^{1/4-9\varepsilon}$.
Assume
\begin{equation}\label{eq:interacting-jet-conditions}
 2\varpi_{\rm sep}^+\le\bar\kappa_1^{N_*},\quad
 N_*\bar\kappa_1^{1-s}\le1/2,\quad
 \varpi_0\le1/4,\quad 8N_*t_2\le t_x .
\end{equation}
There is a finite positive measure $v_{\rm sf}$
on $[-\pi,\pi]^3\times\{\chi_*=1\}$, of
mass in $[1-\varpi_0,1]$, such that
\begin{equation}\label{eq:interacting-small-field-representation}
 \frac{\bar\nu^{a(t)}(\mathfrak G)}{\bar\nu(F_N)}
       =\int \mathrm e^{\mathcal G(t;\theta,A)}\,v_{\rm sf}(\dd\theta\,\dd A).
\end{equation}
The normalized complement
$L(t)=\bar\nu^{a(t)}(F_N\setminus\mathfrak G)/\bar\nu(F_N)$
is entire. Its fourth jet in each variable has
coefficient norm at radius $r$ at most
$\lambda_L^{N_*}$.
\end{proposition}

\begin{proof}
First, the union bound and
\cite[(7)]{TerminalPressure} give
\begin{equation}\label{eq:interacting-small-field-tail}
 \bar\mu(F_N\setminus\mathfrak G^-)\le\varpi_0 .
\end{equation}
For $i=k_*-j\ge1$, write
$v_j=v_*+i(\log L)/2$.
The scale conditions
$v_*\ge8$, $p_0v_*^{p_0-1}\ge28$ give
$p_{0,j}\ge p_{0*}+14i\log L$.
Since $(p_{0,j}-1)^2\ge p_{0,j}$,
the individual tail is at most
$2\mathrm e^{-p_{0*}/2}L^{-7i}$.
Multiply by the coordinate count
$2\ell_*^3L^{3i}$ and sum $L^{-4i}$.
The $3\ell_*^3$ coarse plaquettes give the
second term in \eqref{eq:interacting-complement-constants}.
These estimates yield \eqref{eq:interacting-small-field-tail},
with room in the coefficient $18$.

The exact density identity defines
\begin{equation}\label{eq:interacting-small-field-measure}
 v_{\rm sf}(\dd\theta\,\dd A)=
   \frac{\chi_*(A)Y_*^{\rm x}(A+c_\theta;0)}
             {\bar\nu(F_N)}\,\frac{\dd\theta}{(2\pi)^3}
                                      \dd\mathsf m_*(A).
\end{equation}
We check positivity of this density, which
does not follow from the formal Grassmann
representation. For real marks, the original
fine law $\nu^{\theta,a(t)}$ is nonnegative.
The identity \cite[(21)]{TerminalPressure} says, for every
nonnegative bounded Borel $\Phi$,
\[
 \int\Phi(A)\chi_*(A)Y_*^{\rm x}(A+c_\theta;\mathscr C(t))
                \,\dd\mathsf m_*(A)
       =\nu^{\theta,a(t)}(\mathbf1_{\mathfrak G}\Phi(A_*))\ge0.
\]
The terminal density is continuous. On the
interior of the compact convex set
$\{\chi_*=1\}=\{|dA|_\infty\le p_*\}$, a
nonzero imaginary part or negative real
part would persist on a ball of positive
Gaussian measure, contradicting this identity.
At the boundary, approach along $uA$,
$u\uparrow1$; zero is an interior point.
Thus the density is nonnegative everywhere
on this set. \Cref{prop:interacting-factor}
makes it nonzero there, hence strictly positive.
This reasoning also proves
$\mathcal G(t;\theta,A)\in\R$ for real $t$:
its exponential is a positive ratio and its
continuous value at zero is zero.

Integration of the exact terminal identity over
$\theta$, followed by
\eqref{eq:interacting-terminal-ratio}, proves
\eqref{eq:interacting-small-field-representation}.
At $t=0$ the mass is $\bar\mu(\mathfrak G)$,
which lies in $[1-\varpi_0,1]$.
All integrals here are well-defined Borel
integrals: the density is continuous on
the compact support and zero elsewhere.

For real $t=ez$ with $|t|_\infty\le t_2$,
put $G_z=\sum_cz_cd\widehat{\mathfrak a}_c$.
Disjoint supports imply
$\|G_z\|^2=v_\eta|z|_2^2\le1/4$.
The real source shift gives
\[
 L(ez)=\mathrm e^{-v_\eta|z|_2^2/2}
 \bar\mu\left(\mathbf1_{(F_N\setminus\mathfrak G)+G_z}
                              \mathrm e^{\langle G_z,F^{\rm u}\rangle}\right).
\]
Every photon coordinate changes by at most
$\|G_z\|\le1$ under the shift. Thus the
shifted complement is contained in
$F_N\setminus\mathfrak G^-$.
Cauchy--Schwarz, \eqref{eq:interacting-small-field-tail},
and Gaussian domination at $2G_z$ give
\begin{equation}\label{eq:interacting-real-complement}
 0\le L(t)\le\sqrt{\varpi_0}\,
          \mathrm e^{v_\eta|z|_2^2/2}\le\varpi_L .
\end{equation}
This is the only place where the lowered
thresholds of $\mathfrak G^-$ are needed.

The source-shift identity also gives, for complex marks,
\[
 |\mathring\Phi(t)|\le
       \exp\left(\frac{v_\eta}{2e^2}
                                      \sum_c(\operatorname{Im}t_c)^2\right).
\]
It follows directly by bounding the expectation
of the real exponential and canceling its
real quadratic normalization.
On $|t|_\infty\le t_2$ this is at most
$\mathrm e^{\mathbf D/2}$.
The small-field integral has modulus at most
$2$. Indeed, evenness and the double zero of
$\ell_0$ give $|\ell_0(t)|\le B_0|t|^2/t_x^2$,
while the zero of $w_c$ gives
$|w_c(t)|\le(3/8)|t|/t_x$.
Their sum over cells at
$t_2\le t_x/(8N_*)$ is at most
$1/(256N_*)+3/64$. The remaining logarithm
components sum to at most
$2^{N_*}\varpi_{\rm sep}^+
\le(2\bar\kappa_1)^{N_*}/2\le N_*^{-N_*}/2$,
by \eqref{eq:interacting-jet-conditions}.
Thus $|\mathcal G|\le1/2$ and
$|\mathrm e^{\mathcal G}|\le2$.
Since $L$ is the difference between the
full generating function and the small-field
integral, $|L|\le M_L$ on that polydisc.
It is entire because the finite determinant
density is entire and is bounded on compact
mark sets uniformly in the real photon
field and real twist. The modulus bound
in the external-response section supplies
the required Gaussian integrable majorant
also for either Borel event.

We use the following two-constants estimate.
Let $n\ge1$ be an integer. If a holomorphic function
on a polydisc in $\C^n$ of radius $R$ has modulus at most $M$, and
at most $\delta\le M$ on its real points,
then its modulus is at most $\delta^{1/3}M^{2/3}$
on the polydisc of radius $R/(2n)$.
On the upper
half of the unit disc the harmonic function
$\omega(z)=\frac2\pi\arg[(1+z)/(1-z)]$
is zero on the diameter and one on the
semicircle. Comparing $\log|f(R'z)|$
with $(1-\omega)\log\delta+\omega\log M$,
for $R'<R$, the maximum principle gives
the same inequality in the half disc.
To handle its endpoints, subtract
$\epsilon\log[4/|1-z^2|]$ before applying
the principle, then let $\epsilon\downarrow0$.
Its harmonic conjugate gives a holomorphic
function to which the modulus principle
applies. On $|z|\le q$, $\omega(z)\le
(4/\pi)\arctan q$; reflection treats the
lower half disc. Let $R'\uparrow R$.
Apply this one-variable estimate successively
with $q=1/(2n)$. The remaining exponent of
$\delta$ is at least
$(1-\frac4\pi\arctan(1/(2n)))^n
\ge1-2/\pi>1/3$, proving the assertion.
The case $\delta=0$ follows by the identity
theorem, and $M=0$ is immediate.

Apply it to $L$, with $n=N_*$ and $R=t_2$.
Cauchy's coefficient bound gives
\[
 |L_\alpha|\le\varpi_L^{1/3}M_L^{2/3}(2N_*/t_2)^{|\alpha|}.
\]
In the fourth jet, every exponent is at most
four. Since $y=2N_*r/t_2\ge1$,
\[
 \sum_{\alpha_c\le4}|L_\alpha|r^{|\alpha|}
   \le\varpi_L^{1/3}M_L^{2/3}
                 (1+y+\cdots+y^4)^{N_*}
   \le\lambda_L^{N_*}.
\]
Reality on real marks gives real jet coefficients.
\end{proof}

